% Part: second-order-logic
% Chapter: metatheory
% Section: loewenheim-skolem

\documentclass[../../../include/open-logic-section]{subfiles}

\begin{document}

\olfileid{sol}{met}{lst}

\olsection{فشل مبرهنة لوفنهايم--سكولم في منطق الرتبة الثانية}

\begin{explain}
مؤدى مبرهنة لوفنهايم--سكولم النزولية أن كل مجموعة من
!!{sentence}s لها نموذج لانهائي، لها نموذج !!a{enumerable}.
وهذه أيضًا من ثمرات الاكتمال؛ فبرهانه ينشئ لكل مجموعة متسقة
من !!{sentence}s نموذجًا، ويكون هذا النموذج !!{enumerable}.
وللمبرهنة صورة صعودية كذلك: إذا كان لمجموعة من !!{sentence}s
نموذج !!a{denumerable}، كان لها نموذج !!a{nonenumerable}.
ولا تصح الصورة النزولية ولا الصعودية في منطق الرتبة الثانية.
\end{explain}

\begin{thm}
\ollabel{thm:sol-no-ls} لا تصح مبرهنة لوفنهايم--سكولم في
منطق الرتبة الثانية؛ فبعض !!{sentence}s له نماذج لانهائية،
وليس له نموذج !!{enumerable}.
\end{thm}

\begin{proof}
عرفنا أن
\[
\fn{Count} \ident \lexists[{\OLId{latin-ordinary}{z}}][\lexists[{\OLId{latin-ordinary}{u}}][\lforall[{\OLId{latin-looped}{X}}][(({\OLId{latin-looped}{X}}({\OLId{latin-ordinary}{z}}) \land
      \lforall[{\OLId{latin-ordinary}{x}}][({\OLId{latin-looped}{X}}({\OLId{latin-ordinary}{x}}) \lif {\OLId{latin-looped}{X}}({\OLId{latin-ordinary}{u}}({\OLId{latin-ordinary}{x}})))]) \lif \lforall[{\OLId{latin-ordinary}{x}}][{\OLId{latin-looped}{X}}({\OLId{latin-ordinary}{x}})])]]]
\]
تصدق في !!a{structure}~$\Struct{M}$ إذا وفقط إذا كان
$\Domain{{\OLId{latin-looped}{M}}}$ !!{enumerable}. لذلك تصدق~$\lnot\fn{Count}$
في~$\Struct{M}$ إذا وفقط إذا كان~$\Domain{{\OLId{latin-looped}{M}}}$ !!{nonenumerable}.
ومثل هذه !!{structure}s موجود؛ إذ نأخذ مجموعة !!{nonenumerable}
ونجعلها !!{domain}، كـ$\Pow{\Nat}$ أو~$\Real$.
فتكون لـ$\lnot \fn{Count}$ نماذج لانهائية، ولا يكون لها
نموذج !!{enumerable}.
\end{proof}

\begin{thm}
لبعض !!{sentence}s نماذج !!{denumerable}، ولا نموذج لها
!!{nonenumerable}.
\end{thm}

\begin{proof}
إن~$\fn{Count} \land \fn{Inf}$ تصدق في~$\Nat$، ولا تصدق
في !!{structure}~$\Struct{M}$ مجموعة حملها !!{nonenumerable}.
\end{proof}

\end{document}
